\documentclass[11pt,a4paper]{article}
\usepackage[utf8]{inputenc}
\usepackage[T1]{fontenc}
\usepackage[brazil]{babel}
\usepackage[margin=2.4cm]{geometry}
\usepackage{hyperref}
\usepackage{longtable}
\usepackage{listings}
\title{UniSync --- Integração dos painéis}
\author{Equipe UniSync}
\date{13 de novembro de 2026}
\begin{document}
\maketitle
\tableofcontents
\section{Visão geral}
O sistema mantém o shell Django como entrada de autenticação. O Vue apresenta os painéis, a API .NET autentica e cadastra usuários no banco relacional, e o FastAPI gerencia oportunidades, tarefas e lembretes no MongoDB.
A identidade usada pelo FastAPI vem do JWT emitido pela API .NET.

\section{Autenticação e cadastro}
A entrada local é \url{http://localhost:8081/login/}.
O cadastro de professores fica em \url{http://localhost:8081/cadastro/professor/} e exige o código fornecido pela equipe.
O shell envia nome, e-mail, senha e \texttt{codigoProfessor} para \texttt{/api/usuarios/professores}. A validação do código pertence à API .NET.
Alunos são cadastrados em \texttt{/api/usuarios}. O login usa \texttt{/api/usuarios/login}.
O shell guarda o token na sessão e o fornece ao Vue pelo objeto \texttt{UNISYNC\_SHELL}. As chamadas ao FastAPI usam o cabeçalho Authorization com Bearer.
O cadastro Vue antigo foi substituído por encaminhamento ao shell, eliminando a entrada simulada no dashboard.
O modo \texttt{AUTH\_MOCK} fica desabilitado por padrão e só é ativado explicitamente nos testes.
As requisições de autenticação não seguem redirecionamentos HTTP para outros provedores.
Os nomes de claims contendo schemas.microsoft.com são identificadores do formato JWT e não endpoints de autenticação.
A assinatura, o emissor, o destinatário e a expiração são validados no FastAPI. A leitura de claims no frontend orienta a interface.
A opção Sair encerra a sessão Django.

\section{Permissões}
\begin{longtable}{p{4cm}p{10cm}}
Recurso & Regra \\ \hline
Oportunidades & Consulta pública conforme contrato existente. Apenas professores criam. Atualização e exclusão exigem autoria. Nome e ID do autor vêm do JWT. \\
Lembretes & Alunos e professores autenticados gerenciam somente seus próprios lembretes. O ID do proprietário não é aceito no corpo da requisição. \\
Kanban & Disponível para alunos autenticados. Cada aluno acessa somente suas tarefas, conforme a regra existente do FastAPI. \\
\end{longtable}
Lembretes de outro usuário retornam 404 nas operações de alteração e exclusão. O filtro MongoDB inclui simultaneamente o identificador do documento e o proprietário. Atualizações usam \texttt{\$set} atômico.
As origens locais autorizadas no CORS incluem as portas 5173 e 8081 de localhost e 127.0.0.1.

\section{Lembretes persistentes}
A coleção \texttt{lembretes} contém \texttt{usuario\_id}, \texttt{titulo} e \texttt{concluido}.
O título tem de 1 a 200 caracteres após remoção de espaços nas extremidades.
Há índice por usuário e ordenação de listagem pelo identificador.
\begin{longtable}{p{2cm}p{7cm}p{5cm}}
Método & Rota & Resultado \\ \hline
GET & /lembretes/listar & Lista do usuário \\
POST & /lembretes/criar & Criação, HTTP 201 \\
PATCH & /lembretes/atualizar/\{id\} & Alteração de título ou conclusão \\
DELETE & /lembretes/excluir/\{id\} & Exclusão, HTTP 204 \\
\end{longtable}
O PATCH rejeita corpo vazio, campos desconhecidos e valores nulos. Operações concluídas atualizam a interface após confirmação da API. Erros preservam os dados de entrada.
O painel lateral permite adicionar, editar, concluir, reabrir e excluir lembretes.

\section{Dashboard e oportunidades}
Dashboard e página do professor compartilham \texttt{OportunidadesPanel.vue}, o formulário e a janela modal.
O dashboard inicia com oportunidades futuras. A página do professor inclui também eventos passados.
Ambos limitam a consulta a 15 itens em ordem de data. Ao clicar, o detalhe atualizado é buscado na API e apresenta disciplina, professor, data, tipo e descrição.
Professores autores recebem as ações de editar e excluir. A exclusão solicita confirmação.
Datas sem fuso retornadas pela API são interpretadas como UTC e exibidas no horário local. Formulários convertem o horário local para ISO UTC.
Disciplinas foi removida do menu e sua antiga rota encaminha ao dashboard.

\section{Calendário}
O calendário calcula os dias e o alinhamento de cada mês, incluindo anos bissextos. Permite avançar e retroceder meses, selecionar uma data e voltar para hoje. O dia atual é atualizado periodicamente.
A seleção do calendário é visual e não cria eventos nem altera os filtros das oportunidades.

\section{Kanban}
A tela usa a sessão autenticada e as rotas existentes de tarefas.
Os estados são A\_FAZER, EM\_ANDAMENTO e CONCLUIDO.
SortableJS implementa arraste com mouse e toque. O DOM é restaurado antes da aplicação da resposta para manter o Vue responsável pelo estado.
Arrastar para a lixeira abre uma confirmação. Também há controles de status e exclusão por botão para teclado.
A edição altera o título. Durante uma requisição, os controles ficam bloqueados para evitar operações simultâneas conflitantes. Em falha, o estado confirmado é preservado.
As animações respeitam a preferência por movimento reduzido.

\section{Organização do código}
\begin{itemize}
\item FastAPI: modelos, schemas, repositório e router próprios para lembretes.
\item Vue: serviços HTTP em services, estado do Kanban em composables, componentes reutilizáveis em components.
\item Os utilitários de sessão, datas e URLs de autenticação concentram regras compartilhadas.
\item Os testes Python cobrem validação, autenticação, autoria e CORS. Os testes frontend cobrem calendário, sessão, datas e fluxos de navegador.
\end{itemize}

\section{Execução local}
Na raiz do projeto, após alterações de dependências ou do backend, executar:
\begin{verbatim}
docker compose up -d --build --no-deps \
  api-python frontend-web frontend-shell
\end{verbatim}
No PowerShell, executar o comando acima em uma única linha, sem as barras de continuação.
O Compose de desenvolvimento monta os arquivos locais do shell e Vue. Alterações nas dependências e no FastAPI exigem reconstrução da imagem.
O endereço de uso é \url{http://localhost:8081/app/}.
Variáveis relevantes: API\_BASE\_URL no shell, VITE\_PYTHON\_API\_URL e VITE\_SHELL\_URL no frontend, e as configurações MongoDB/JWT existentes no FastAPI.
O acesso ao site da equipe fica no menu lateral em \href{https://teste2234234.my.canva.site/in-cio}{Conheça a equipe}.

\section{Testes e limites de validação}
Nesta revisão passaram 78 testes do FastAPI sem E2E, 22 testes do shell Django, 3 testes unitários JavaScript e 6 testes de navegador. A compilação do frontend e a análise ESLint dos arquivos alterados passaram. Há um aviso de depreciação do TestClient/httpx no ambiente Python. A validação do CRUD com login real continua sendo uma etapa de integração do ambiente.
\begin{verbatim}
# backend-python
python -m pytest -m "not e2e"
# frontend-shell
python manage.py test authentication
# frontend-web-aluno
npm run test:unit
npx playwright install chromium
npm run test:browser
npm run build
\end{verbatim}
Os testes de navegador interceptam a API com dados controlados para validar interações e erros sem modificar o banco usado pela equipe. A simulação preserva os itens durante recargas da página.
A validação completa em ambiente real requer uma conta de aluno, uma conta de professor com código válido e MongoDB de teste. Conferir criação, recarga, alteração, conclusão e exclusão dos lembretes, e arraste e recarga das tarefas. Usar duas contas para confirmar isolamento.
Não armazenar senhas, tokens ou códigos reais neste documento ou nos testes.
\end{document}

